\section{Condiciones holon\'omicas del selector gravitatorio general}

La segunda plantilla gravitatoria parte de
\begin{equation}
\label{rm:eq:gravity-delta-theta}
\delta_\theta
=\frac{A_{\rm deg}-C^*_{\rm deg}}{360}
\end{equation}
y de una longitud de ruta
\(\Lambda_{27}=\mathcal L_{\rm hol}(\widetilde{\mathcal K}_{27})\)
en un levantamiento enriquecido del retorno. La expresi\'on
\begin{equation}
\label{rm:eq:gravity-holonomic-template}
\boxed{
G_{\rm hol}
=\frac{c_{\rm int}^3}{\hbar_{\rm int}}
\bigl(\delta_\theta\Lambda_{27}\bigr)^2}
\end{equation}
es dimensionalmente exacta. Su promoci\'on a selector general requiere que
\(\mathcal L_{\rm hol}\) sea natural bajo cambio c\'iclico de base,
conjugaci\'on, reetiquetado y subdivisi\'on, y que posea una normalizaci\'on
\'unica estable bajo refinamiento.

Hay adem\'as un control de tors\'ion: todo car\'acter multiplicativo positivo
\(\chi:K_{27}\to\R_+^\times\) sobre un ciclo observable finito de orden
cuatro es trivial, porque \(\chi(g)^4=1\) y \(\chi(g)>0\) implican
\(\chi(g)=1\). Una respuesta holon\'omica no trivial debe, por tanto,
definirse sobre el levantamiento no torsional o sobre la memoria
enriquecida; no puede
obtenerse renombrando un car\'acter positivo del ciclo finito.

La carta
\begin{equation}
\label{rm:eq:gravity-decimal-chart}
10^{-11}\left(6+\frac{674}{1000}+\frac{30}{10^5}\right)
=6.67430\times10^{-11}
\end{equation}
es una evaluaci\'on exacta de entradas declaradas. No fija
\(\Lambda_{27}\) ni demuestra que \cref{rm:eq:gravity-holonomic-template}
adopte ese valor.

\begin{keyresult}
La condici\'on de periodicidad determina \(G\) sin emplear su expansi\'on
decimal: la estructura cronol\'ogica CT108 selecciona \((54/\pi)^2\), la
recta gravitatoria fija el tipo dimensional y la familia bidireccional
conserva \(G\hbar\). El acoplamiento resultante se representa sobre el
portador tensorial construido por \(P_5\), mientras que la reducci\'on
f\'isica transversal sin traza se efect\'ua mediante un mapa distinto de
rango dos.
\end{keyresult}

\begin{falsifier}
El bloque circular queda refutado si
\(r_g(\theta)\neq\theta(\pi/54)\ell_0\), si existe otra soluci\'on positiva
de \cref{rm:eq:gravity-closing-condition}, o si
\(G_a\hbar_a/c^3\neq\theta_{108}^{\rm circ}\ell_0^2\). El bloque tensorial queda
refutado si \(\pi_{\rm ico}\) no preserva incidencia o no entrelaza la
acci\'on declarada, si su estabilizador no es \(C_5\), si \(P_5\) no es
idempotente autoadjunto de traza cinco, si
\(\Gamma_0\Gamma_0^*\neq(8/5)I\), o si \(\Lambda(n)\) no tiene rango dos.
La lectura de cinco polarizaciones queda refutada si los tensores de
\cref{rm:eq:gravity-helicity2-real,rm:eq:gravity-helicity1-real,rm:eq:gravity-helicity0-real}
no son una base ortonormal o si \(\Lambda(n)\) retiene un peso distinto de
\(\pm2\).
El selector holon\'omico general queda resuelto por obstrucci\'on: sin una
\(\mathcal L_{\rm hol}\) natural, normalizada y estable no existe una
selecci\'on un\'ivoca compatible con todas las relabelizaciones; usar
\cref{rm:eq:gravity-decimal-chart} para elegirla incurre en circularidad e
invalida la construcci\'on.
\end{falsifier}

\begin{statusbox}
Transductor, unicidad del selector bajo la condici\'on de periodicidad CT108 y
familia de Planck: \status{Teorema interno exacto bajo la condici\'on geom\'etrica declarada}.
El descenso \(\Omega_{12}^{\rm ico}\simeq A_5/C_5\) es exacto bajo las cuatro
hip\'otesis expl\'icitas de \(\pi_{\rm ico}\); la construcci\'on HMT de ese
mapa pertenece al corredor excepcional de procedencia. Fijado el descenso,
proyector \(P_5\), entrelazador \(\Gamma_5\), base de cinco helicidades,
reducci\'on TT y espectro de \(\mathbb G_5\): \status{Teorema interno exacto}.
Lectura masiva, ecuaciones de campo y realizaci\'on
Palatini--Cartan--Holst: \status{Realizaci\'on f\'isica tipada}. Plantilla
holon\'omica: \status{Demostrada en homogeneidad}; la inexistencia de un
selector general un\'ivoco sin normalizaci\'on adicional es una
\status{obstrucci\'on de canonicidad demostrada}. Carta SI:
\status{Contraste metrol\'ogico externo}.
\end{statusbox}
